\appendix
\section{Consolidated Release-Wide Operations Reference}
\label{app:release-operations}

This appendix follows a run from installation through solver selection,
initial data, execution, output, and analysis. Use the derivations and module
chapters for the equations and update sequence. The initial-condition, I/O,
parallel, GPU, and visualization chapters supply the detailed operating
instructions.

\subsection{Installation and Toolchain Requirements}
\label{sec:installation}

VENUS builds in its source checkout using dependencies installed by the
user. Keep the sibling
\path{cpu} and \path{gpu} directories together: the GPU build consumes the
CPU unequal-Plummer calibration table and its public field-clock headers
from \path{cpu/include}.

\subsubsection{CPU prerequisites}

The CPU core requires a C++17 compiler and GNU-compatible \code{make}.
The portable targets use \code{std::thread}; the performance targets require
an OpenMP runtime.  MPI is optional and is required only for targets ending in
\code{_mpi}.  Python~3.8 or newer is used by IC generators, analysis,
and figures.
The release Python requirements are NumPy, SciPy, Matplotlib, Pillow, ImageIO,
pypdf (for vector-curve digitization).  A clean environment can
be prepared with
\par\noindent\begin{minipage}{\linewidth}
\codesource{Release environment and CPU build verification}{\path{requirements.txt}, \path{Makefile}, and \path{cpu/Makefile}}{Creates the documented Python environment, builds the portable CPU task front ends, and exposes their help interfaces}
\begin{lstlisting}[language=bash]
cd VENUS
python3 -m venv .venv
. .venv/bin/activate
python3 -m pip install -r requirements.txt
make cpu
\end{lstlisting}
\end{minipage}
On Apple Clang, use the serial targets unless an OpenMP runtime was installed
explicitly.  On clusters, load a consistent compiler/MPI stack. The MPI wrapper and C++
runtime must belong to compatible builds, because the wrapper links the
executable to those libraries.

\subsubsection{CUDA prerequisites}

The GPU tree requires CMake~3.18 or newer, a CUDA toolkit, a host compiler
supported by that toolkit, and Python~3.8 or newer for build-time generation of the
subhalo table.  CUDA-aware MPI is not required by the correctness-first
multi-GPU layer: marker migrations and radial moment reductions are host
staged.  An MPI C++ development package is required when
\code{VENUS_GPU_ENABLE_MPI=ON}.  Select the actual compute capability at
configuration time; \code{80} is appropriate for NVIDIA A100 hardware:
\par\noindent\begin{minipage}{\linewidth}
\codesource{Single-GPU configure and build}{\path{gpu/CMakeLists.txt}}{Configures the CUDA architecture and CPU calibration-table dependency, then builds the GPU libraries and task executables}
\begin{lstlisting}[language=bash]
cmake -S gpu -B build/gpu \
  -DCMAKE_BUILD_TYPE=Release \
  -DCMAKE_CUDA_ARCHITECTURES=80 \
  -DVENUS_CPU_REFERENCE_INCLUDE="$PWD/cpu/include"
cmake --build build/gpu -j 8
\end{lstlisting}
\end{minipage}
The GPU build reads the CPU calibration table so both implementations use
the same Plummer encounter data. Configuration therefore requires that table
at its documented source location.

\subsubsection{Release-wide targets}

The root \path{Makefile} is an orchestration layer; CPU numerical flags remain
in \path{cpu/Makefile} and GPU flags remain in \path{gpu/CMakeLists.txt}.
\par\noindent\begin{minipage}{\linewidth}
\codesource{Release-wide build and verification entry points}{\path{Makefile}}{Delegates each named command to the owning CPU, GPU or manual build system without changing its numerical contract}
\begin{lstlisting}[language=bash]
make help
make cpu                 # serial reference executables
make cpu-openmp          # OpenMP executables
make cpu-mpi             # supported MPI executables
make cpu-tools           # compile every maintained C++ runner/diagnostic
make cpu-check-setup
make gpu                 # configure and build CUDA tree
make gpu-mpi             # configure/build the MPI CUDA variant
make manual
\end{lstlisting}
\end{minipage}
For a non-default compute capability, pass the override to the final target,
for example \code{make gpu GPU_ARCH=90} or
\code{make gpu-mpi GPU_ARCH=90}.  Each dependent target runs configuration;
an architecture supplied only to an earlier standalone
\code{gpu-configure} invocation is not retained as a Make variable by a later
\code{make} process.
Build directories and run output are not source.  They remain outside the
implementation subtrees under \path{build/} and \path{output/}.

\subsection{Choosing a Solver}
\label{sec:solver-choice}

Choose the physical equation first and the representation second:
\begin{enumerate}
\item Use grid Einstein--Vlasov when the relativistic metric and a resolved
 phase-space distribution are essential.  Use the EV WPIC path only for its
stated low- or high-compactness task; it has no collision operator.
\item Use grid VP when a deterministic distribution, explicit phase-space
flux, Hamiltonian interpolation, or conservative SLDG/canonical analysis is
required.
\item Use CPU WPIC when long characteristic orbits, adaptive field/collision
events, positive unequal marker masses, or subhalo scattering/capture are the
natural representation and marker-number convergence can be measured.
\item Use CPU SphDSMC for the persistent-marker SIDM benchmark family and
when direct spherical shell pairing, orbit retry, and mean-free-path control
are desired.
\item Use GPU WPIC when the supported Newtonian marker physics dominates the
work and the documented restart, background, boundary, device, and MPI
contracts match the calculation.  Do not select it for EV,
SphDSMC, canonical-orbit, or grid semi-Lagrangian calculations.
\end{enumerate}

For collisional work, distinguish the physical targets.  Material SIDM treats
markers as sampling an underlying dark-matter material and uses a microscopic
cross section per mass.  The subhalo module treats each marker as a
complete extended Plummer object; it uses calibrated scattering and capture
cross sections, unequal object masses, internal energy, and merger
compaction.  A constant material cross section, a velocity power law, a
Yukawa/Born differential kernel, and a calibrated subhalo table are not
interchangeable ways to set the same scalar parameter.

\subsection{Configuration, Command Lines, and Reproducible Runs}
\label{sec:configuration-shared}

The CPU modular solver uses an INI-like file with command-line/environment
overrides as described in the CPU configuration reference.  CPU WPIC runners
also expose purpose-specific command-line controls.  The GPU task executables read a weighted-PIC restart and command-line
options. To move a calculation between CPU and GPU, identify the same physical
model, particle state, and numerical controls in their respective interfaces.

Every run should preserve, next to its output:
\begin{itemize}
\item the exact input file or complete command line;
\item source identity and dirty-tree status where available;
\item executable/build identity, compiler, CUDA architecture, and loaded
cluster modules;
\item thread count, MPI rank count, and rank-to-GPU binding;
\item RNG seed and persistent event epoch/counters;
\item initial restart hash and final restart hash;
\item a normal-completion, walltime-checkpoint, locality-stop, or other
unambiguous terminal-status record.
\end{itemize}
Start a fresh run in an empty output directory. Continue from an atomic
checkpoint using a distinct continuation directory or the selected runner's
documented resume mechanism. This keeps the input checkpoint and earlier
output available when comparing the two segments.

\subsubsection{Parameter changes that alter the physical problem}

Grid bounds, background expansion, Poisson source convention, cross-section
law, object-size/mass relation, inner/outer boundary conditions, and capture
enablement define the physical model.  Time tolerances, mesh counts, marker
counts, cloud order, collision cell occupancy, majorant safety, and output
cadence define a discretization or algorithmic realization.  Record both, but
do not present agreement after changing the first group as a resolution test.

\subsection{Initial Conditions and CPU--GPU Handoff}
\label{sec:initial-data-shared}

Grid EV/VP/BP runs accept analytic distributions or the documented
\code{EVICRWL1} binary format.  Specialized Python generators construct
Eddington-inverted halo equilibria, infalling shells, Frank cell-average
states, and persistent SphDSMC marker populations.  Each generator states its
mass normalization and radial/velocity/angular-momentum convention; a binary
reader must verify dimensions, monotonic grids, finiteness, and exact payload
length before allocating the evolution state.

WPIC and SphDSMC initial states are not interchangeable.  A WPIC marker owns
represented mass and orbit coordinates; a SphDSMC marker additionally
participates in its solver's persistent collision-shell and restart contract.
For the GPU, generate or import a supported CPU WPIC checkpoint, validate it
with the CPU inspector, and then read it through the versioned GPU restart
loader.  The GPU does not contain an independent equilibrium sampler in this
release.  This prevents nominally identical CPU/GPU comparisons from starting
from different Monte Carlo populations.

For any sampled IC, archive the seed, marker count, sampling proposal, weight
rule, rejected-sample count, normalization correction, and binned residual
against the target density/velocity moments.  Successful generation establishes a readable sampled state. Equilibrium
requires a collisionless evolution check of density, velocity moments, and
conserved quantities.

\subsection{Output and Restart Contract}
\label{sec:io-shared}

CPU grid solvers write the CSV profile/history, phase projections, metadata,
and restart products detailed in the CPU I/O section.  Marker runners add
effective marker number, collision histories, pair/capture ledgers, locality,
and event-clock fields.  GPU output names are listed in the GPU chapter.  Code
that consumes output must locate columns by header name rather than by an
unversioned numeric position.

The minimum interpretation rules are:
\begin{itemize}
\item \code{time} is physical solver time; a file index is not a time.
\item Density is a physical volume density under the solver's mass and radial
normalization; deposited shell mass is not density until divided by volume.
\item Radial and tangential second moments must be combined according to the
header and coordinate convention; \(L^2/r^2\) supplies tangential speed.
\item An integrated collision age and a minimum local collision time answer
different questions and must not share a plot label.
\item Escaped, hidden/delayed, captured, internal, and active masses remain
separate ledger terms.  Closing a total by silently renormalizing active mass
destroys the diagnostic.
\end{itemize}

A restart preserves the state required for the next numerical update.
It must retain the population/distribution, mesh contract, time and step,
field and collision event ages, random epoch/counters, delayed-boundary state,
internal energy, unique/source IDs, and every model option whose omission
would change subsequent evolution.  Readers check magic, version, endian
sentinel, counts, finite fields, enumeration ranges, monotonic grids, and
trailing/truncated bytes.  Restart equivalence compares an uninterrupted run
against a checkpoint/resume run at a common endpoint.

\subsection{Plotting, Analysis, and GIF Construction}
\label{sec:plotting-shared}

The maintained catalogue is \path{notes/SCRIPT_REFERENCE.md}.  Plotting
follows a four-stage contract:
\begin{enumerate}
\item discover inputs from one explicit run root and reject mixed schemas;
\item read header-named columns, validate finiteness/monotonic radius, and
associate every curve with physical time;
\item compute derived density slopes, dispersions, core measures, or
collision coverage under one documented smoothing and radial mask;
\item write the figure/GIF plus sidecars recording source count, selected
times, last rendered curve, normalization, command/provenance, and completion.
\end{enumerate}

Representative commands are:
\par\noindent\begin{minipage}{\linewidth}
\codesource{EV and SIDM analysis and visualization tools}{\path{cpu/tools/plot_einstein_vlasov_outputs.py}, \path{cpu/tools/render_density_slope_sigma2_gif.py}, \path{cpu/tools/summarize_venus_collision_history.py}, and \path{gpu/scripts/render_venus_gpu_density_slope_gif.py}}{Plots Einstein--Vlasov output, renders spherical evolution, summarizes collision histories, and renders the GPU density/slope sequence}
\begin{lstlisting}[language=bash]
python3 cpu/tools/plot_einstein_vlasov_outputs.py \
  --input-dir output/cpu/ev-run --output-dir output/cpu/ev-run/figures
python3 cpu/tools/render_density_slope_sigma2_gif.py \
  --input-dir output/cpu/run --output output/cpu/run/evolution.gif
python3 cpu/tools/summarize_venus_collision_history.py \
  output/cpu/run/collision_history.csv
python3 gpu/scripts/render_venus_gpu_density_slope_gif.py --help
\end{lstlisting}
\end{minipage}
The GPU front end accepts explicit collisionless/collisional run directories,
output directory, output name, and profile-selection controls.  Use each
script's \code{--help}; the exact option names are authoritative.
GIF frames are temporary derived products.  They use an explicit work
directory or the operating-system temporary directory, are deleted only after
a successful assembly, and never rely on a hard-coded user's
\path{/private/tmp}.  Animation timing must come from the frame-time sidecar or
physical-time selection, not from lexicographic filename spacing.

\subsection{Parallel Execution and Accelerator Semantics}
\label{sec:parallel-shared}

OpenMP and MPI change ownership of work, not the continuum equation.  The CPU
thread wrapper assigns independent output cells or shells to workers and uses
thread-local reductions.  Supported grid EV/VP MPI builds decompose the
angular-momentum axis into contiguous slabs and perform the global moment
reductions required by the metric/Poisson solve.  Collision paths that lack
distributed ownership reject a multi-rank launch before evolution; silently
running an independent collision system on every rank is forbidden.

The GPU assigns markers or radial cells to CUDA threads and collision shells
to blocks.  Atomic moment deposition combines concurrent contributions while preserving
the represented sums to floating-point tolerance. The order of additions can
vary, which changes their last binary digits.  Multi-GPU VP reduces
small radial moment arrays globally while each rank retains local markers.
The supported subhalo path migrates complete physical collision cells so a
cell is owned by one rank.  Global RNG coordinates use seed, epoch, physical
cell, substep, trial, and draw dimension rather than local thread identity.

For performance reports, include initialization and I/O policy, marker/cell
counts, steps or physical interval, collision strength, accepted/proposed
events, hardware model, rank/GPU mapping, thread/block configuration, and
separate kernel/transfer/output timings.  A faster run with fewer collision
proposals or a different timestep is not a speedup measurement of the same
problem.

\subsection{Failure Modes and Safe Diagnosis}
\label{sec:failure-modes-shared}

Common fail-closed conditions and their interpretation are:
\begin{itemize}
\item occupied phase space reaches a grid boundary: enlarge or redesign the
IC/domain and repeat a resolution study; do not mask the diagnostic;
\item compactness gap or metric floor is crossed: preserve the last valid
state and interpret only under the EV stopping contract;
\item collision probability/majorant exceeds its bound: reduce the collision
substep or fix the majorant; clipping an invalid probability biases rates;
\item \(\Delta r/\lambda_{\rm mfp}>1\): refine physical collision cells and
marker sampling; decreasing output cadence is irrelevant;
\item positivity or conservation residual exceeds tolerance: isolate
transport, field, and collision operators before interpreting a quantitative
run;
\item MPI rank count is unsupported: use one rank plus threads, or the
explicitly supported MPI module;
\item CUDA reports no device: submit through the GPU scheduler partition; a
successful host compile is not a device test;
\item restart reading fails: retain both files and diagnose header/state
differences; never repair a binary checkpoint by editing bytes;
\item plotter reports missing/mixed columns: point it at one compatible run
and use the renderer for that named-column schema; do not fill absent physical
quantities with zeros.
\end{itemize}

\section{Numerical Verification}
\label{app:developer-verification}

Implementation verification covers structural interfaces, parsers, restarts,
invariants and CPU--GPU comparisons. Developer regression suites are
maintained separately from the distribution. The distributed build targets
compile the production solvers and analysis tools.

A scientific calculation also requires the problem-specific convergence
programme described in Section~\ref{sec:failure-modes-shared}. Report the
spatial, velocity, angular-momentum, marker-number, timestep, collision-shell,
box-size and seed comparisons that determine the accuracy of its physical
observables.
